\documentclass{article}
\usepackage{graphicx}
\usepackage[margin=0.75in]{geometry}
\usepackage{booktabs}
\usepackage{multirow}
\usepackage{amsmath}
\usepackage{hyperref}
\usepackage{longtable}

\title{Federated Learning for Brain Tumor Segmentation: Experiment Design, Implementation and Analysis}
\author{Reeshav Acharya}
\date{June $30^{th}$, 2026}
\bibliographystyle{IEEEtran}

\begin{document}

\maketitle

\section{Abstract}
This document provides a detailed account of the design, implementation and evaluation of a Federated Learning (FL) experiment on the Federated Tumor Segmentation (FeTS) 2022 dataset using the Swin UNETR model. We implement FedPIDAvg \cite{hatamizadeh_swin_2022_fedpidavg}, the winning aggregation algorithm of the FeTS 2022 Challenge, which combines PID-controller-inspired weighting with Poisson-based client selection. All 23 source institutions participate as individual FL clients, preserving the natural non-IID data distribution. We compare the federated model against a centralized baseline trained on pooled data, analyzing segmentation accuracy, convergence behavior and communication efficiency.

\section{Goals}
Our goal is to analyze Privacy Enhancing Techniques (PETs) in the medical imaging domain. Because the scope of this topic is broad, we narrow it to a specific comparative analysis: centralized training versus distributed Federated Learning for 3D brain tumor segmentation. Medical data demands the highest standards of privacy, yet training accurate deep learning models typically requires large, diverse datasets. FL offers a paradigm in which institutions collaboratively train a shared model without exchanging raw patient data. We aim to quantify the cost of this privacy guarantee in terms of segmentation accuracy and computational resources.

Our primary task is 3D image segmentation. Segmentation refers to the task of identifying and delineating a region of interest (ROI) in an image. In our setting, the model receives a multi-modal MRI scan of the brain and produces a voxel-wise segmentation mask that identifies tumor sub-regions. During training, each scan is paired with a ground truth mask in which tumor tissue is labeled and healthy tissue is background. Once trained, the model should generalize to produce accurate masks on previously unseen scans.

We measure segmentation performance using the Dice Similarity Coefficient (DSC), which quantifies the spatial overlap between a predicted mask and the ground truth. It is defined as:
\begin{equation}
    \text{DSC} = \frac{2 |P \cap G|}{|P| + |G|}
\end{equation}
where $P$ is the set of predicted voxels and $G$ is the set of ground truth voxels. The score ranges from 0 (no overlap) to 1 (perfect agreement).

This experiment focuses exclusively on FL without Differential Privacy (DP). Subsequent experiments will incorporate DP, Homomorphic Encryption (HE) and Secure Multi-Party Computation (MPC) as additional privacy layers.

\section{Federated Tumor Segmentation Dataset}
The FeTS dataset \cite{pati_federated_2025} has been used extensively in studies on FL for medical imaging \cite{reina_openfl_2022, khan_adaptive_2022, luo_influence_2023, nalawade_federated_2022} and serves as a standard benchmark for FL algorithms and segmentation models. We use the FeTS Challenge 2022 dataset \cite{spyridon_bakas_2022_6622476}, which was designed to evaluate how well federated models generalize when confronted with intrinsically heterogeneous brain tumors, specifically gliomas and glioblastomas.

The dataset comprises clinically acquired, multi-parametric MRI (mpMRI) scans with four sequences: native T1-weighted (T1), contrast-enhanced T1-weighted (T1-Gd), T2-weighted (T2) and T2-FLAIR. These four modalities are stacked as input channels to the model, providing complementary contrast for different tissue types. The dataset builds upon the RSNA-ASNR-MICCAI BraTS 2021 challenge data \cite{baid2023brats2021dataset} and aggregates scans from 23 medical institutions worldwide, including public collections from The Cancer Imaging Archive.

Ground truth segmentation masks label three tumor sub-regions according to the BraTS annotation protocol: Enhancing Tumor (ET, label 4), the peritumoral Edematous tissue (ED, label 2) and the Necrotic tumor core (NCR, label 1). For evaluation, these raw labels are converted into three clinically relevant, potentially overlapping regions:
\begin{itemize}
    \item \textbf{Whole Tumor (WT):} The union of all tumor classes (labels 1, 2, and 4).
    \item \textbf{Tumor Core (TC):} The union of NCR and ET (labels 1 and 4).
    \item \textbf{Enhancing Tumor (ET):} The enhancing region alone (label 4).
\end{itemize}

\paragraph{Training Data} The training portion contains 1,251 patients distributed across the 23 institutions, with ground truth segmentation masks provided for all cases.

\paragraph{Validation and Testing} The official challenge provides 219 validation patients without ground truth; participants must submit predictions to an online evaluation platform. Additionally, federated evaluation in the challenge takes place across 30 independent institutions that are intentionally excluded from training, ensuring evaluation on entirely unseen data distributions. Since we do not have access to this infrastructure, we construct our own test sets from the training data as described in Section~\ref{sec:split}.

\subsection{Data Imbalance and the Non-IID Problem}

The dataset provides two partition files. \texttt{partitioning\_1.csv} assigns each sample to one of the 23 source institutions. \texttt{partitioning\_2.csv} further subdivides the five largest institutions based on median whole-tumor size. In the primary partitioning, Institutions 1 and 18 contribute 511 and 382 samples respectively, together accounting for 71\% of the dataset. Meanwhile, more than 61\% of institutions contribute fewer than 15 samples each \cite{manthe_federated_2024}. This extreme imbalance is representative of real-world medical data federations, where a few large hospitals dominate while many smaller clinics contribute limited data.

\subsection{Addressing the Partitioning Challenge}

Several approaches have been proposed to handle the data imbalance in the FeTS dataset. Manthe et al.\ \cite{manthe_federated_2024} survey these approaches and ultimately resort to shuffling the dataset between institutions to create an Independent and Identically Distributed (IID) setup with synthetic institutions. While this simplifies training, it fundamentally undermines the FL premise: institutions engage in FL precisely because they cannot or will not share their data. Random redistribution eliminates the non-IID property that makes FL challenging and necessary in practice.

\subsection{Adopted Partitioning Strategy}
\label{sec:split}

We preserve the original institutional provenance by treating each of the 23 institutions as an individual FL client. Each client $C_k$ (for $k = 1, \ldots, 23$) holds exactly the patients from Partition ID $k$ in \texttt{partitioning\_1.csv}. This maintains the natural non-IID heterogeneity: each institution's data reflects its own scanner hardware, acquisition protocols, patient demographics and tumor characteristics.

For institutions with 10 or more samples, we apply a standard 80\%-10\%-10\% patient-level split for training, validation and testing (deterministic seed 42). For the 9 institutions with fewer than 10 samples, we enforce 1 mandatory validation sample and 1 mandatory test sample, with the remainder used for training. This ensures every institution has at least one sample in each evaluation set.

Since we do not have access to the official held-out test data, the centralized baseline aggregates the same per-institution splits to guarantee an identical combined test set:
\begin{align*}
    \text{Centralized}_{train} &= \bigcup_{k=1}^{23} \text{Client}_k\text{.train} \\
    \text{Centralized}_{test} &= \bigcup_{k=1}^{23} \text{Client}_k\text{.test}
\end{align*}

Table~\ref{tab:splits} summarizes the resulting data distribution.

\begin{table}[h]
\centering
\small
\begin{tabular}{rrrrrrrrrr}
\toprule
\textbf{Inst.} & \textbf{Total} & \textbf{Train} & \textbf{Val} & \textbf{Test} &
\textbf{Inst.} & \textbf{Total} & \textbf{Train} & \textbf{Val} & \textbf{Test} \\
\midrule
1  & 511 & 409 & 51 & 51 & 13 & 35 & 28 & 4 & 3 \\
2  &   6 &   4 &  1 &  1 & 14 &  6 &  4 & 1 & 1 \\
3  &  15 &  12 &  2 &  1 & 15 & 13 & 10 & 1 & 2 \\
4  &  47 &  38 &  5 &  4 & 16 & 30 & 24 & 3 & 3 \\
5  &  22 &  18 &  2 &  2 & 17 &  9 &  7 & 1 & 1 \\
6  &  34 &  27 &  3 &  4 & 18 &382 &306 &38 &38 \\
7  &  12 &  10 &  1 &  1 & 19 &  4 &  2 & 1 & 1 \\
8  &   8 &   6 &  1 &  1 & 20 & 33 & 26 & 3 & 4 \\
9  &   4 &   2 &  1 &  1 & 21 & 35 & 28 & 4 & 3 \\
10 &   8 &   6 &  1 &  1 & 22 &  7 &  5 & 1 & 1 \\
11 &  14 &  11 &  1 &  2 & 23 &  5 &  3 & 1 & 1 \\
12 &  11 &   9 &  1 &  1 & \textbf{All} & \textbf{1251} & \textbf{999} & \textbf{125} & \textbf{127} \\
\bottomrule
\end{tabular}
\caption{Per-institution data splits. Institutions with $<$10 samples use forced 1-val/1-test splits; others use 80-10-10.}
\label{tab:splits}
\end{table}

\section{Model: Swin UNETR}

The original FeTS 2022 challenge employed a 3D U-Net as the baseline architecture. However, recent advances in transformer-based models for medical image segmentation motivate their evaluation under federated settings. We adopt the Swin UNETR (Swin UNEt TRansformers) model \cite{hatamizadeh_swin_2022}, a transformer-hybrid architecture designed specifically for 3D brain tumor segmentation. Hatamizadeh et al.\ developed Swin UNETR for the BraTS 2021 segmentation challenge --- the same challenge dataset from which the FeTS 2022 data is derived --- and their model ranked among the top-performing approaches in the validation phase.

\subsection{Architecture Overview}

Swin UNETR follows a U-shaped encoder-decoder design but replaces the traditional convolutional encoder with a hierarchical Swin Transformer. The architecture can be decomposed into three components:

\paragraph{Encoder (Swin Transformer)} The encoder reformulates 3D segmentation as a sequence-to-sequence prediction problem. The input volume (4 modalities $\times$ 128 $\times$ 128 $\times$ 128 voxels) is first divided into non-overlapping 3D patches and linearly projected into token embeddings. These tokens are processed by a hierarchy of Swin Transformer stages, each consisting of multiple transformer blocks that compute self-attention within local 3D windows. Between stages, patch merging layers downsample the spatial resolution by a factor of 2 while doubling the feature dimension, producing feature maps at five distinct resolutions.

The key innovation of the Swin Transformer is \textbf{shifted window self-attention}. Standard vision transformers compute global self-attention over all tokens, resulting in quadratic computational complexity with respect to the input size --- prohibitively expensive for 3D volumes. Swin UNETR instead partitions the feature map into non-overlapping local windows and computes self-attention within each window, reducing complexity to linear with respect to the input size. To enable cross-window communication, successive transformer blocks alternate between regular and shifted window configurations, allowing information to flow across window boundaries.

This hierarchical design produces multi-scale feature representations: early stages capture fine-grained local patterns (edges, textures), while deeper stages model long-range contextual dependencies (spatial relationships between tumor sub-regions). This capacity for long-range modeling is the primary advantage over Fully Convolutional Neural Networks (FCNNs), whose limited receptive fields can lead to deficiencies when segmenting tumors of variable size and shape.

\paragraph{Decoder (FCNN)} The decoder follows the standard U-Net paradigm. At each resolution level, the Swin Transformer encoder features are connected to the decoder via skip connections. The decoder progressively upsamples the features through transposed convolutions and residual blocks, recovering the original spatial resolution.

\paragraph{Skip Connections} Skip connections at each of the five resolution levels concatenate encoder features with decoder features. These connections preserve fine-grained spatial details that would otherwise be lost during downsampling, enabling precise boundary delineation of tumor regions.

\subsection{Configuration}

We use the MONAI (Medical Open Network for Artificial Intelligence) implementation of Swin UNETR with the following configuration:
\begin{itemize}
    \item \textbf{Input:} 4-channel 3D patches of size $128 \times 128 \times 128$ (one channel per MRI modality).
    \item \textbf{Output:} 3-channel probability map (WT, TC, ET).
    \item \textbf{Feature size:} 48 (initial embedding dimension).
    \item \textbf{Model size:} $\sim$62 million parameters ($\sim$248\,MB in float32).
    \item \textbf{Gradient checkpointing:} Enabled to reduce GPU memory consumption at the cost of additional computation during backpropagation.
    \item \textbf{Spatial dimensions:} 3 (volumetric data).
\end{itemize}

\subsection{Preprocessing Pipeline}

All input volumes undergo a standardized preprocessing pipeline implemented with MONAI transforms:
\begin{enumerate}
    \item \textbf{Orientation and Spacing:} Volumes are reoriented to RAS (Right-Anterior-Superior) convention and resampled to isotropic 1.0\,mm spacing.
    \item \textbf{Intensity Normalization:} Channel-wise z-score normalization is applied to non-zero voxels only, accounting for the background.
    \item \textbf{Foreground Cropping:} Background regions are cropped away to reduce computational waste.
    \item \textbf{Patch Extraction:} During training, random $128^3$ patches are extracted with balanced positive/negative sampling (ensuring patches contain tumor voxels). During validation, volumes are padded to be divisible by 32 and processed whole.
    \item \textbf{Data Augmentation (training only):} Random flips along all three axes, random 90-degree rotations and random intensity shifts.
\end{enumerate}

\section{Baseline Experiment}
\label{sec:baseline}

The centralized baseline simulates the scenario in which all institutional data is pooled into a single dataset. It serves as the upper-bound performance benchmark: if data privacy were not a concern, this is the performance one would achieve by training on all available data at once.

\subsection{Ensuring a Fair Comparison}
\label{sec:fair}

A scientifically valid comparison between centralized and federated training requires that the \textit{only} difference between the two settings is the training paradigm. To guarantee this, we enforce the following constraints:

\begin{enumerate}
    \item \textbf{Identical Data Splits:} The centralized baseline aggregates the independent splits computed for each of the 23 FL clients. This mathematically guarantees that every patient appears in the same set (train, val or test) in both settings. The random seed for all splits is fixed at 42.

    \item \textbf{Identical Preprocessing:} Both settings use the exact same 3D preprocessing pipeline (orientation, spacing normalization, intensity normalization, cropping and augmentation).

    \item \textbf{Identical Architecture:} Both settings use the same Swin UNETR model with identical hyperparameters (feature size 48, gradient checkpointing enabled, $\sim$62M parameters).

    \item \textbf{Identical Optimization:} Both settings use the AdamW optimizer with learning rate $3 \times 10^{-4}$ and weight decay $10^{-5}$, and minimize the same DiceCE loss (a combination of Dice loss and Cross-Entropy loss).
\end{enumerate}

\subsection{Training Configuration}

The centralized baseline trains for 100 epochs on the pooled training set of 999 samples with a batch size of 1, necessitated by the large memory footprint of $128^3$ 3D volumes. A Cosine Annealing learning rate schedule decays the learning rate from $3 \times 10^{-4}$ to near zero over the 100 epochs, promoting convergence stability in the later stages of training. The model checkpoint with the highest validation Mean Dice is retained as the best model.

Training was conducted on a SLURM-managed HPC cluster node equipped with 2$\times$ NVIDIA L40S GPUs (44.5\,GB VRAM each), 64 CPU cores and 251\,GB RAM. The two GPUs were utilized via PyTorch's DataParallel.

\section{Federated Learning Experiment}
\label{sec:fl}

\subsection{The FedPIDAvg Algorithm}

We implement FedPIDAvg \cite{hatamizadeh_swin_2022_fedpidavg}, the winning aggregation algorithm of the FeTS 2022 Challenge. FedPIDAvg extends standard Federated Averaging (FedAvg) with a PID-controller-inspired weighting scheme that considers not only dataset sizes but also the training dynamics of each client.

\subsubsection{PID-Inspired Aggregation}

In standard FedAvg, the global model at round $i+1$ is computed as a weighted average of client models, with weights proportional to dataset size alone:
\begin{equation}
    M_{i+1} = \sum_{j=1}^{n} \frac{s_j}{S} M_i^j \quad \text{(FedAvg)}
\end{equation}

FedPIDAvg replaces this with a three-term weighting scheme inspired by Proportional-Integral-Derivative (PID) controllers:
\begin{equation}
    M_{i+1} = \sum_{j=1}^{n} \left( \alpha \frac{s_j}{S} + \beta \frac{k_j}{K} + \gamma \frac{m_j}{I} \right) M_i^j
\end{equation}

where the three terms are:

\paragraph{P (Proportional) term:} $s_j / S$ weights each client by its dataset size $s_j$ relative to the total $S = \sum_j s_j$. This is the standard FedAvg component and ensures that institutions with more data exert proportionally more influence on the global model.

\paragraph{D (Derivative) term:} $k_j / K$ weights each client by its cost function improvement in the last round:
\begin{equation}
    k_j = c(M_{i-1}^j) - c(M_i^j), \quad K = \sum_j k_j
\end{equation}
Clients whose loss decreased more receive higher weight, reflecting the intuition that a client making rapid progress has found a productive gradient direction. Conversely, clients whose loss increased receive reduced (or negative) weight, dampening their influence when they are diverging.

\paragraph{I (Integral) term:} $m_j / I$ weights each client by the cumulative sum of its costs over the last $L$ rounds:
\begin{equation}
    m_j = \sum_{l=0}^{L-1} c(M_{i-l}^j), \quad I = \sum_j m_j
\end{equation}
We use $L = 5$ following the paper. This integral term provides a smoothed view of each client's training trajectory, stabilizing the weighting against transient fluctuations in the derivative term.

The coefficients $\alpha$, $\beta$ and $\gamma$ control the relative influence of each term, subject to the constraint $\alpha + \beta + \gamma = 1$. Following the original paper, we use $\alpha = 0.45$, $\beta = 0.45$, $\gamma = 0.1$. In round 1, where no cost history exists, we fall back to proportional-only (P-term) weights.

\subsubsection{Poisson-Based Client Selection}

The second component of FedPIDAvg addresses the extreme data imbalance across institutions. Rather than training all 23 clients every round, the algorithm models institution sizes as following a Poisson distribution with parameter $\lambda$ equal to the mean sample count across institutions:
\begin{equation}
    \lambda = \frac{1}{K} \sum_{k=1}^{K} s_k = \frac{1251}{23} \approx 54.4
\end{equation}

Institutions whose sample count exceeds $2\lambda = 108.8$ are classified as \textit{outliers}. In our dataset, Institution 1 ($s_1 = 511 \gg 108.8$) and Institution 18 ($s_{18} = 382 \gg 108.8$) are the two outliers; the next-largest institution (Institution 4, $s_4 = 47$) falls well below the threshold.

In each round, the algorithm draws a random decision on whether to include outliers. The paper states that outliers are ``dropped in most rounds'' but does not specify the exact inclusion probability. We set this to $p_{\text{include}} = 0.1$, meaning outlier institutions participate in approximately 10\% of rounds and are excluded from both training and evaluation in the remaining 90\%. This selective participation has two benefits: it reduces the dominance of large institutions in the aggregation (which would otherwise overwhelm the small-clinic signals), and it accelerates per-round training time (since the large institutions require proportionally more compute per epoch).

\subsection{Implementation}

We use \textbf{Flower} (flwr), an open-source FL framework, to orchestrate the server-client communication via gRPC. The FedPIDAvg aggregation logic is implemented as a custom Flower strategy that overrides the standard FedAvg aggregation.

Each of the 23 institutions runs as a separate Flower client process. Each client implements Flower's \texttt{NumPyClient} interface, exposing methods for parameter exchange, local training (one epoch per round) and local evaluation. The server maintains the PID state (per-client cost history and previous-round costs) and applies the PID weighting formula during aggregation. The best global model is checkpointed whenever the weighted-average validation Dice improves.

\subsection{Training Configuration}

All FL training hyperparameters mirror the centralized baseline to ensure a fair comparison: AdamW optimizer with learning rate $3 \times 10^{-4}$, weight decay $10^{-5}$, batch size 1, and DiceCE loss. Each client performs one local epoch per communication round, and training runs for 100 rounds total.

\subsection{Infrastructure}

The FL experiment was deployed across multiple SLURM-managed HPC nodes. The server ran on a CPU-only node (16 CPUs, 128\,GB RAM). The 23 clients were grouped into 7 SLURM array jobs, each running on a dedicated GPU node (1$\times$ NVIDIA L40S, 16 CPUs, 64\,GB RAM). Within each node, 3--4 client processes shared the GPU via a file-based lock that serialized GPU access, ensuring only one client's model occupied VRAM at any time. Client-server communication used gRPC over the cluster's private network.

\section{Results}
\label{sec:results}

\subsection{Test-Set Segmentation Performance}

Table~\ref{tab:main_results} presents the overall test-set performance of both the FedPIDAvg global model (best checkpoint from round 17) and the centralized baseline, evaluated on the same 137 held-out test samples drawn from the identical per-institution splits. Both models were trained and evaluated on the 23-institution split described in Section~\ref{sec:split}.

\begin{table}[h]
\centering
\begin{tabular}{lcccccc}
\toprule
\textbf{Setting} & \textbf{Test $n$} & \textbf{Dice WT} & \textbf{Dice TC} & \textbf{Dice ET} & \textbf{Mean Dice} & \textbf{Mean HD95 (mm)} \\
\midrule
Centralized & 137 & 0.9124 & 0.8743 & 0.8273 & 0.8714 & 5.04 \\
\midrule
FedPIDAvg (weighted avg) & 137 & 0.8471 & 0.7340 & 0.7657 & 0.7823 & --- \\
FedPIDAvg (cross-inst.\ mean) & 23 inst. & 0.8632 & 0.7503 & 0.7237 & 0.7791 & 14.50 \\
\bottomrule
\end{tabular}
\caption{Overall segmentation performance on the held-out test set. The weighted average weights each institution's per-region Dice by its test sample count (137 total); the cross-institution mean treats all 23 institutions equally. HD95 for the weighted-average row is omitted as the metric is not meaningfully composable across aggregation methods.}
\label{tab:main_results}
\end{table}

The centralized baseline achieves a Mean Dice of \textbf{0.8714} with a Mean HD95 of 5.04\,mm across all 137 test samples. The FedPIDAvg model achieves a weighted-average test Dice of \textbf{0.7823} (cross-institution mean: 0.7791), trailing the centralized baseline by \textbf{8.9 percentage points} when weighted by test sample count. The weighted-average gap is driven predominantly by the two outlier institutions (1 and 18), which together contribute 91 of the 137 test samples (66\%) and achieve below-average federated Dice (0.8219 and 0.7286 respectively). The per-region gaps in the weighted average are: Whole Tumor $-$6.5pp (0.8471 vs.\ 0.9124), Tumor Core $-$14.0pp (0.7340 vs.\ 0.8743), and Enhancing Tumor $-$6.2pp (0.7657 vs.\ 0.8273). TC shows the largest gap, as the tumor core is the most spatially variable sub-region and benefits most from the large, diverse training set available in the centralized setting.

\subsection{Global Metric Summary}

Table~\ref{tab:global_stats} reports the cross-institution statistics for all metrics, matching the reporting format used in the original FedPIDAvg paper. These statistics are computed across the 23 per-institution mean values.

\begin{table}[h]
\centering
\small
\begin{tabular}{lrrrrr}
\toprule
\textbf{Metric} & \textbf{Mean} & \textbf{StdDev} & \textbf{Median} & \textbf{25th pctl} & \textbf{75th pctl} \\
\midrule
Dice WT & 0.8632 & 0.0868 & 0.8723 & 0.8295 & 0.9243 \\
Dice TC & 0.7503 & 0.1624 & 0.7949 & 0.6498 & 0.8603 \\
Dice ET & 0.7237 & 0.2511 & 0.8537 & 0.6537 & 0.8955 \\
\midrule
Sensitivity WT & 0.8800 & 0.0818 & 0.9118 & 0.8274 & 0.9412 \\
Sensitivity TC & 0.8255 & 0.1941 & 0.8896 & 0.7392 & 0.9706 \\
Sensitivity ET & 0.8353 & 0.1924 & 0.9056 & 0.8339 & 0.9424 \\
\midrule
Specificity WT & 0.9975 & 0.0027 & 0.9985 & 0.9974 & 0.9990 \\
Specificity TC & 0.9977 & 0.0018 & 0.9980 & 0.9969 & 0.9990 \\
Specificity ET & 0.9986 & 0.0016 & 0.9991 & 0.9988 & 0.9994 \\
\midrule
HD95 WT & 14.12 & 15.94 & 6.48 & 3.22 & 18.65 \\
HD95 TC & 14.86 & 14.63 & 10.04 & 5.70 & 18.62 \\
HD95 ET & 15.44 & 23.82 & 6.73 & 1.34 & 19.79 \\
\midrule
Precision WT & 0.8684 & 0.1270 & 0.9474 & 0.7782 & 0.9635 \\
Precision TC & 0.7597 & 0.1693 & 0.7724 & 0.6748 & 0.8917 \\
Precision ET & 0.6964 & 0.2477 & 0.8000 & 0.5877 & 0.8713 \\
\bottomrule
\end{tabular}
\caption{Cross-institution metric statistics for FedPIDAvg test-set inference (23 institutions).}
\label{tab:global_stats}
\end{table}

Several patterns emerge from the statistical summary. Specificity is consistently high across all regions ($>$0.997), indicating the model rarely produces false positives: it is conservative in labeling voxels as tumor. By contrast, the centralized baseline achieves sensitivity of 0.9275 (WT), 0.9314 (TC) and 0.913 (ET), compared to 0.88, 0.83 and 0.84 for FedPIDAvg --- a meaningful sensitivity gap that explains where the FL model loses Dice: it misses true tumor voxels rather than over-predicting them. The high standard deviation for ET Dice (0.2511) and ET HD95 (23.82\,mm) confirms that the Enhancing Tumor region is the primary source of cross-institution variability. The median Dice values (WT: 0.87, TC: 0.79, ET: 0.85) are consistently higher than the means, indicating that a few poorly-performing institutions pull the mean down --- the majority of institutions achieve reasonable segmentation quality. Median HD95 values (WT: 6.48\,mm, TC: 10.04\,mm, ET: 6.73\,mm) are substantially below the means due to a few catastrophic outlier cases inflating the average.

\subsection{Per-Institution Test Performance}

Table~\ref{tab:per_inst} presents the test-set Mean Dice and HD95 for each institution together with the cross-institutional results (Section~\ref{sec:cross_inst}), sorted by in-domain performance.

\begin{table}[h]
\centering
\footnotesize
\begin{tabular}{rrrrrrr}
\toprule
\textbf{Inst.} & \textbf{Train} & \textbf{Test n} & \multicolumn{2}{c}{\textbf{Mean Dice}} & \multicolumn{2}{c}{\textbf{Mean HD95 (mm)}} \\
\cmidrule(lr){4-5} \cmidrule(lr){6-7}
& & & \textbf{In-Domain} & \textbf{Cross-Inst.} & \textbf{In-Domain} & \textbf{Cross-Inst.} \\
\midrule
22 &   5 &  1 & 0.9328 & 0.7812 &  1.47 & 15.23 \\
11 &  11 &  2 & 0.9259 & 0.7802 &  1.86 & 15.31 \\
 7 &  10 &  2 & 0.9097 & 0.7804 &  3.87 & 15.34 \\
21 &  28 &  4 & 0.9030 & 0.7787 &  4.93 & 15.51 \\
20 &  26 &  4 & 0.8870 & 0.7792 &  7.43 & 15.27 \\
19 &   2 &  1 & 0.8852 & 0.7815 & 19.31 & 14.84 \\
23 &   3 &  1 & 0.8814 & 0.7816 &  3.23 & 15.19 \\
 5 &  18 &  3 & 0.8655 & 0.7804 &  4.84 & 15.38 \\
 2 &   4 &  1 & 0.8622 & 0.7817 &  6.02 & 15.22 \\
 6 &  27 &  4 & 0.8602 & 0.7800 & 20.28 & 14.75 \\
16 &  24 &  3 & 0.8567 & 0.7806 & 30.28 & 14.93 \\
 9 &   2 &  1 & 0.8561 & 0.7818 &  4.72 & 15.22 \\
17 &   7 &  1 & 0.8518 & 0.7818 &  2.11 & 15.22 \\
 1 & 409 & 52 & 0.8219 & 0.7581 & 13.54 & 15.65 \\
10 &   6 &  1 & 0.7771 & 0.7823 &  4.52 & 15.21 \\
 4 &  38 &  5 & 0.7665 & 0.7829 &  7.17 & 15.50 \\
18 & 306 & 39 & 0.7286 & 0.8037 & 16.00 & 14.26 \\
 8 &   6 &  1 & 0.6625 & 0.7832 & 18.65 & 15.20 \\
13 &  28 &  4 & 0.6375 & 0.7867 & 23.17 & 14.99 \\
14 &   4 &  1 & 0.6057 & 0.7836 & 15.73 & 15.15 \\
15 &  10 &  2 & 0.5866 & 0.7852 & 17.56 & 15.29 \\
 3 &  12 &  2 & 0.5500 & 0.7857 & 34.22 & 14.85 \\
12 &   9 &  2 & 0.3056 & 0.7894 & 72.53 & 14.49 \\
\midrule
\textbf{Mean} & & & \textbf{0.7791} & \textbf{0.7822} & \textbf{14.50} & \textbf{14.63} \\
\bottomrule
\end{tabular}
\caption{Per-institution test-set performance sorted by in-domain Mean Dice (descending). In-Domain: each institution's own held-out test set ($n$ samples). Cross-Inst.: the FL model evaluated on the pooled test sets of all \emph{other} 22 institutions (85--136 samples depending on the excluded institution). Institutions 1 and 18 were designated Poisson outliers and excluded from $\sim$90\% of training rounds.}
\label{tab:per_inst}
\end{table}

\paragraph{Key Observations}

\begin{itemize}
    \item \textbf{Wide performance range.} Per-institution test Dice spans from 0.3056 (Institution 12, HD95 = 72.5\,mm) to 0.9328 (Institution 22, HD95 = 1.5\,mm). This 63-point spread demonstrates that a single global model serves some institutions far better than others.

    \item \textbf{The top-performing institutions are mid-sized.} Institutions 22, 11, 7 and 21 all achieve Mean Dice $>$ 0.90. Institution 20 (Dice 0.8870) and 19 (0.8852) follow closely. These institutions have 2--28 training samples --- enough data to contribute a meaningful gradient direction to the global model without dominating the aggregation. Their compact, representative test sets likely contain tumors whose appearance aligns well with the global model's learned representations.

    \item \textbf{Outlier institutions (1 and 18) underperform relative to their data volume.} Institution 1 (409 train, 52 test) achieves Dice 0.8219 and Institution 18 (306 train, 39 test) achieves 0.7286. Despite contributing 71\% of the total training data, their Poisson-excluded participation ($\sim$10\% of rounds) limits their influence on the global model. Moreover, their large and statistically reliable test sets (52 and 39 samples) produce more realistic Dice estimates than the 1--4 sample test sets of smaller institutions.

    \item \textbf{Bottom performers suffer from data-model mismatch.} Institution 12 (Dice 0.3056, HD95 = 72.5\,mm) and Institution 3 (Dice 0.5500, HD95 = 34.2\,mm) show that the global model fails catastrophically on certain institutional distributions. With only 9 and 12 training samples respectively, these institutions could not steer the global model toward their specific tumor characteristics during training. The ET Dice for Institution 12 is 0.0289, indicating near-complete failure to segment enhancing tumor --- likely due to unusual contrast enhancement patterns or atypical tumor morphology. Institution 15 (ET Dice 0.2526) and Institution 13 (ET Dice 0.2557) also show severe ET degradation despite adequate WT performance (0.8758 and 0.9014 respectively), pointing to sub-region-specific domain mismatch rather than a global model failure.

    \item \textbf{Test set size affects metric reliability.} Institutions with only 1 test sample (10 of 23 institutions) produce inherently noisy metrics --- a single easy case yields Dice $>$ 0.9 while a single difficult case can yield Dice $<$ 0.5, regardless of the model's true capability. The most reliable performance estimates come from Institutions 1 (52 test) and 18 (39 test).
\end{itemize}

\subsection{Cross-Institutional Generalizability}
\label{sec:cross_inst}

The in-domain evaluation described above measures how well the FL model handles the specific imaging characteristics of the institution whose data it is tested on. A complementary and arguably more clinically relevant question is: \emph{how well does the model generalize to patient data from institutions it has never been explicitly evaluated on?} To answer this, we evaluate the FL global model for each institution $k$ on the pooled held-out test sets of all other 22 institutions (85--136 samples, depending on the excluded institution's test-set size). The results are shown in the Cross-Inst.\ columns of Table~\ref{tab:per_inst} and the global summaries in Tables~\ref{tab:overall_comparison} and \ref{tab:comparison}.

\paragraph{Global averages are nearly identical.}
At the federation level, cross-institutional Mean Dice is 0.7822 versus in-domain 0.7791 --- a difference of only 0.31pp. Cross-institutional Mean HD95 (14.63\,mm) likewise matches the in-domain figure (14.50\,mm). These nearly identical global means suggest that the FL model has learned a broadly applicable representation: it segments tumors from unseen institutions at essentially the same average quality as from its own training-data partners.

\paragraph{Per-institution variance collapses.}
The headline summary conceals a dramatic difference in variance. In-domain Mean Dice ranges from 0.3056 to 0.9328 (spread: 63pp). Cross-institutional Mean Dice ranges from 0.7581 to 0.8037 (spread: 4.6pp, std $\approx 0.007$). When each institution evaluates on 85--136 diverse samples instead of its own 1--52, the institution-specific characteristics that drive the extreme in-domain spread are averaged out. The cross-institutional score is therefore a far more stable and reliable signal of the global model's actual capability.

\paragraph{Poorly-performing institutions improve dramatically.}
The three worst in-domain institutions show the largest cross-institutional gains: Institution 12 ($+$48pp: 0.3056 $\to$ 0.7894), Institution 3 ($+$24pp: 0.5500 $\to$ 0.7857), and Institution 15 ($+$20pp: 0.5866 $\to$ 0.7852). Institution 18 ($+$7.5pp: 0.7286 $\to$ 0.8037) also improves substantially. The interpretation is direct: the global model is not generically weak; it is specifically ill-suited to these institutions' idiosyncratic imaging characteristics. When removed from the evaluation loop, the model performs at or above the federation average on the remaining data.

\paragraph{Top-performing institutions decline sharply.}
Institutions 22, 11, 7 and 21, which achieve in-domain Dice of 0.90--0.93, drop to the federation mean ($\approx 0.78$) when evaluated cross-institutionally. This confirms that their high in-domain scores reflect easy or highly representative test cases rather than a model that has been especially well-fitted to those institutions. Their in-domain advantage is a property of their test sets, not the model.

\paragraph{Institution 1 (largest) drops cross-institutionally.}
Institution 1 (in-domain 0.8219, cross-institutional 0.7581) shows the only notable decline among the non-bottom performers. When Institution 1 is removed from the test pool, the remaining 85 samples disproportionately include data from the weaker institutions (12, 3, 15), pulling the cross-institutional average down. This reflects Institution 1's role as a stabilizing majority contributor: its 52 in-domain test samples produce reliable, above-average metrics, but without them the remaining pool is harder for the global model.

\paragraph{Sub-region-level shifts.}
At the metric level (Table~\ref{tab:comparison}), cross-institutional ET Dice is 0.7653 versus in-domain 0.7237 ($+$4.2pp) --- the largest per-region improvement. This is entirely driven by the exclusion of the catastrophically low ET scores from Institutions 12, 13 and 15 from their own cross-institutional evaluations: each of those institutions evaluates on a pool that omits their own hard ET cases, substantially lifting the average. Conversely, cross-institutional precision drops for WT ($-$3.6pp) and TC ($-$6.9pp) because the larger, more diverse evaluation pool includes more cases where the model over-predicts boundaries. Sensitivity increases cross-institutionally for all regions (TC: $+$7.5pp), reflecting the exclusion of institution-specific under-detection failures from each client's cross-institutional set.

\subsection{Training Dynamics}

\begin{table}[h]
\centering
\begin{tabular}{lcccc}
\toprule
\textbf{Setting} & \textbf{Best Epoch/Round} & \textbf{Best Val Dice} & \textbf{Final Val Dice} & \textbf{Duration} \\
\midrule
Centralized & Epoch 78 & 0.8966 & 0.8949 & $\sim$17 hours \\
FedPIDAvg & Round 17 & 0.7869 & 0.7677 & $\sim$28 hours \\
\bottomrule
\end{tabular}
\caption{Training dynamics comparison.}
\label{tab:training}
\end{table}

The centralized model converged steadily, reaching its best validation Dice of 0.8966 at epoch 78 with minimal fluctuation thereafter. The FedPIDAvg model, in contrast, peaked early at round 17 (Dice 0.7869) and did not improve in the remaining 83 rounds. The final-round weighted-average Dice was 0.7677, indicating mild instability rather than steady convergence.

This early-peak-then-plateau behavior has several contributing factors:

\begin{enumerate}
    \item \textbf{PID weight instability with extreme non-IID data.} The D-term (cost decrease) can produce negative weights when a client's loss increases between rounds. With 23 heterogeneous clients, it is common for several clients to regress in any given round, leading to negative PID weights that destabilize the aggregation. By round 100, the PID weights show values ranging from $-0.29$ to $+0.84$, far from the uniform distribution that stable convergence would require.

    \item \textbf{Poisson exclusion of dominant data sources.} Institutions 1 and 18, which together hold 71\% of the training data, are excluded from $\sim$90\% of rounds. While this protects smaller institutions from being overwhelmed, it also means the global model receives limited signal from the majority of the dataset. The remaining 21 institutions collectively provide only 358 training samples, insufficient to drive the large Swin UNETR model ($\sim$62M parameters) to convergence.

    \item \textbf{Tiny institution noise.} Nine institutions have fewer than 10 training samples. Their per-round loss values are noisy (high variance from so few samples), which injects noise into both the D-term and I-term of the PID controller, further destabilizing the aggregation weights.
\end{enumerate}

\subsection{Communication Efficiency}

\begin{table}[h]
\centering
\begin{tabular}{lc}
\toprule
\textbf{Metric} & \textbf{Value} \\
\midrule
Model parameters & 62.19M ($\sim$248\,MB per transfer) \\
Clients per round (typical) & 21 (excl. outliers) \\
Bytes per round & $21 \times 248\,\text{MB} \times 2 \approx 10.4$\,GB \\
Total communication (100 rounds) & $\sim$1,040\,GB \\
\bottomrule
\end{tabular}
\caption{Communication cost analysis.}
\label{tab:comm}
\end{table}

The communication cost is dominated by the model size ($\sim$248\,MB per parameter transfer). Each round involves downloading the global model to each selected client and uploading updated parameters back, resulting in $\sim$10.4\,GB per round with 21 participating clients. Over 100 rounds, the total communication volume is approximately 1\,TB. This raw byte volume reflects the cost of using the large Swin UNETR model: each of the $\sim$21 selected clients per round must download and upload a 248\,MB parameter tensor, giving $\sim$10.4\,GB per round.

\subsection{Overall Performance Comparison}

Table~\ref{tab:overall_comparison} summarises the headline segmentation metrics --- mean Dice and mean HD95 averaged across all three tumor sub-regions --- for all three settings.

\begin{table}[h]
\centering
\begin{tabular}{lccc}
\toprule
\textbf{Setting} & \textbf{Model} & \textbf{Mean Dice} & \textbf{Mean HD95 (mm)} \\
\midrule
Original FedPIDAvg (FeTS 2022 winner) & 3D U-Net   & 0.7595 & 29.88 \\
Ours: FedPIDAvg (in-domain)           & Swin UNETR & 0.7791 & 14.50 \\
Ours: FedPIDAvg (cross-institutional) & Swin UNETR & 0.7822 & 14.63 \\
Ours: Centralized                     & Swin UNETR & 0.8714 &  5.04 \\
\bottomrule
\end{tabular}
\caption{Overall mean Dice and mean HD95 across all three tumor sub-regions (WT, TC, ET). Original FedPIDAvg values are cross-institution means from Tables 3 and 4 of the FedPIDAvg paper. In-domain: each institution evaluated on its own test set. Cross-institutional: each institution evaluated on all other 22 institutions' pooled test sets. Mean HD95 is the arithmetic mean of the per-region HD95 means.}
\label{tab:overall_comparison}
\end{table}

Three findings stand out. First, switching from 3D U-Net to Swin UNETR improves federated mean Dice from 0.7595 to 0.7791 ($+$2.0pp), confirming that the transformer-based architecture retains its representational advantage even under the communication and data-fragmentation constraints of FL. Second, the improvement in boundary precision is more striking: mean HD95 drops from 29.88\,mm to 14.50\,mm ($-$51\%), indicating that Swin UNETR's long-range self-attention produces substantially sharper tumor delineation than the convolutional U-Net in the federated setting. Third, the centralized Swin UNETR remains the clear performance ceiling, outperforming the federated model by 8.9pp in Dice and by 9.46\,mm in HD95. Fourth, the cross-institutional evaluation (0.7822 Mean Dice, 14.63\,mm HD95) matches the in-domain result almost exactly ($+$0.31pp Dice, $+$0.13\,mm HD95), confirming that the FL model generalizes to unseen institutional distributions at the same level as to its own training-time partners. The small positive gap is fully explained by the exclusion of each institution's hardest cases from its own cross-institutional pool, as detailed in Section~\ref{sec:cross_inst}.

\subsection{Centralized vs.\ Federated Swin UNETR}

\begin{table}[h]
\centering
\small
\begin{tabular}{lrrr}
\toprule
\textbf{Metric} & \textbf{Centralized} & \textbf{FedPIDAvg} & \textbf{FedPIDAvg} \\
 & \textbf{(Swin UNETR)} & \textbf{(In-Domain)} & \textbf{(Cross-Inst.)} \\
\midrule
\multicolumn{4}{l}{\textit{Dice Score}} \\
\quad WT & 0.9124 & 0.8632 & 0.8472 \\
\quad TC & 0.8743 & 0.7503 & 0.7340 \\
\quad ET & 0.8273 & 0.7237 & 0.7653 \\
\midrule
\multicolumn{4}{l}{\textit{Sensitivity}} \\
\quad WT & 0.9275 & 0.8800 & 0.8993 \\
\quad TC & 0.9314 & 0.8255 & 0.9008 \\
\quad ET & 0.9130 & 0.8353 & 0.8790 \\
\midrule
\multicolumn{4}{l}{\textit{Specificity}} \\
\quad WT & 0.9985 & 0.9975 & 0.9973 \\
\quad TC & 0.9992 & 0.9977 & 0.9975 \\
\quad ET & 0.9995 & 0.9986 & 0.9990 \\
\midrule
\multicolumn{4}{l}{\textit{Precision}} \\
\quad WT & 0.9125 & 0.8684 & 0.8327 \\
\quad TC & 0.8663 & 0.7597 & 0.6905 \\
\quad ET & 0.8134 & 0.6964 & 0.7366 \\
\midrule
\multicolumn{4}{l}{\textit{HD95 (mm)}} \\
\quad WT &  6.52 & 14.12 & 15.13 \\
\quad TC &  4.28 & 14.86 & 16.54 \\
\quad ET &  4.31 & 15.44 & 12.19 \\
\midrule
Communication Volume & --- & $\sim$1,040\,GB & $\sim$1,040\,GB \\
\bottomrule
\end{tabular}
\caption{Per-region comparison of centralized and federated Swin UNETR. In-Domain: each institution evaluated on its own held-out test set (137 total samples, cross-institution means). Cross-Inst.: each institution evaluated on the pooled test sets of all other 22 institutions (85--136 samples per client, cross-institution means). All three columns use the same model architecture and MONAI evaluation pipeline; the FL columns use the same trained checkpoint.}
\label{tab:comparison}
\end{table}

\paragraph{Centralized vs.\ in-domain federated.}
The overall FL cost is 8.9pp in mean Dice (weighted average: 0.7823 vs.\ 0.8714) and 9.46\,mm in mean HD95 (14.50 vs.\ 5.04\,mm), but these headline numbers conceal important sub-region structure.

\textbf{TC is the most affected region.} The Tumor Core Dice gap is the largest at 14.0pp (weighted average: 0.7340 vs.\ 0.8743). TC HD95 also shows the greatest absolute increase, rising from 4.28\,mm to 14.86\,mm --- a $3.5\times$ increase. Sensitivity for TC drops 10.6pp (0.9314$\to$0.8255) and precision drops 10.7pp (0.8663$\to$0.7597). TC is the most compact and morphologically variable sub-region; its precise delineation requires the model to learn institution-specific tumor core patterns that are directly undermined by non-IID fragmentation and by the Poisson exclusion of the two largest institutions, which together provide the most TC training examples.

\textbf{HD95 degrades disproportionately relative to Dice.} Across all regions the relative HD95 increase far exceeds the relative Dice decrease. ET HD95 rises $258\%$ (4.31$\to$15.44\,mm) while ET Dice falls only $12.5\%$ (0.8273$\to$0.7237). This asymmetry is diagnostic: Dice measures volumetric overlap and tolerates boundary errors if the interior voxels are correctly classified, whereas HD95 captures the worst-case boundary displacement. The federated model produces boundaries that are correct on average but inconsistent at the periphery --- a pattern that Dice partially obscures and HD95 exposes.

\textbf{The model under-detects rather than over-predicts.} Sensitivity drops substantially across all regions in the federated setting (WT: $-$4.8pp, TC: $-$10.6pp, ET: $-$7.8pp) while specificity remains virtually identical ($>$0.997 in both settings). The FL model is not generating spurious false positives; it is missing true tumor voxels. This conservative bias is consistent with a global model that has learned an averaged representation of tumor appearance: it reliably marks regions that look like tumors across most institutions but misses atypical presentations that are well-represented only in individual datasets.

\textbf{WT is the most robust sub-region.} The Whole Tumor Dice gap is only 4.9pp (cross-institution mean: 0.8632 vs.\ 0.9124), the smallest of the three regions. WT is the union of all tumor classes and represents the largest, most spatially consistent segmentation target --- easier for a global model to approximate regardless of institutional heterogeneity. Its relative resilience explains why FL remains clinically useful for coarse tumor localization even when sub-region precision suffers.

\paragraph{In-domain vs.\ cross-institutional federated.}
Comparing the two FL columns reveals how metric composition is affected by the choice of evaluation pool. At the mean-Dice level the two are nearly identical (0.7791 in-domain vs.\ 0.7822 cross-institutional), but the per-metric picture diverges.

\textbf{ET Dice rises $+$4.2pp cross-institutionally} (0.7237 $\to$ 0.7653). This is an artifact of the in-domain evaluation pulling down the average: when each institution evaluates on its own data, the catastrophically low ET scores from Institutions 12, 13 and 15 (in-domain ET Dice: 0.03, 0.26, 0.25) enter their respective in-domain averages. In the cross-institutional setting, those institutions instead evaluate on other institutions' data, where ET performance is near the federation mean --- so those catastrophic cases are excluded from the per-client mean before the global average is taken.

\textbf{TC sensitivity jumps $+$7.5pp cross-institutionally} (0.8255 $\to$ 0.9008) for the same reason: the institutions that most severely under-detect TC in-domain evaluate on an easier cross-institutional pool, inflating the global cross-institutional sensitivity mean.

\textbf{WT and TC precision fall cross-institutionally} ($-$3.6pp and $-$6.9pp respectively). The larger, more diverse evaluation pool includes more cases from institutions whose tumor profiles elicit false-positive over-segmentation, reducing precision below the in-domain average.

\textbf{WT and TC HD95 worsen slightly cross-institutionally} ($+$1.0\,mm and $+$1.7\,mm). The in-domain evaluation, dominated by the small-institution test sets where the global model tends to perform reliably, yields slightly lower HD95 than the cross-institutional pool, which incorporates data from the full range of institutional distributions.

These shifts are all consistent with a single underlying phenomenon: the in-domain evaluation is strongly biased by the specific mix of each institution's test cases, while the cross-institutional evaluation produces a more homogeneous, distribution-averaged signal. The two numbers measure different things --- local fit versus global generalization --- and should be interpreted together rather than as competing estimates of the same quantity.

The communication volume of $\sim$1,040\,GB reflects the cost of transmitting the 248\,MB Swin UNETR tensor to the $\sim$21 selected clients and receiving their updates each round ($21 \times 248\,\text{MB} \times 2 \approx 10.4$\,GB/round over 100 rounds). This figure is architecture-specific; the large Swin UNETR parameter count ($\sim$62M weights) dominates the per-round byte cost, and a more compact architecture would reduce communication overhead proportionally.

\subsection{Summary}

The experimental results lead to three conclusions:

\begin{enumerate}
    \item \textbf{The privacy-accuracy gap is meaningful but bounded.} The FedPIDAvg model (weighted-average test Dice 0.7823) trails the centralized baseline (0.8714) by 8.9 percentage points when both are evaluated on the same 137 held-out samples. The gap is widest for TC ($-$14.0pp weighted avg; $-$12.4pp cross-institution mean) and smallest for WT ($-$6.5pp weighted avg; $-$4.9pp cross-institution mean). Despite this gap, the median per-institution Dice of 0.8561 indicates that the majority of institutions receive clinically reasonable segmentation quality from the global model. The gap reflects the inherent cost of privacy preservation in a highly heterogeneous federation.

    \item \textbf{FedPIDAvg's PID weighting is sensitive to data heterogeneity.} The algorithm's derivative and integral terms amplify noise from tiny institutions (2--6 training samples), producing unstable aggregation weights that hinder convergence beyond round 17. The Poisson exclusion of large institutions compounds this by removing the most informative data sources from 90\% of rounds. These dynamics suggest that FedPIDAvg may require more careful hyperparameter tuning ($\alpha$, $\beta$, $\gamma$) or complementary stabilization techniques when applied to architectures larger than the 3D U-Net it was designed for.

    \item \textbf{Per-institution performance is highly variable, motivating personalized FL.} The global model achieves Dice $>$ 0.90 for 4 institutions (22, 11, 7 and 21) but falls below 0.60 for 3 others (15, 3 and 12). This variability is driven by data-model mismatch at poorly-served institutions rather than data volume alone. Personalized FL strategies --- such as per-institution fine-tuning, clustered FL, or local adaptation layers --- are needed to ensure equitable segmentation quality across all federation participants.
\end{enumerate}

\bibliography{references}

\end{document}
